\documentclass[a4paper,11pt]{article}
\usepackage[utf8]{inputenc}
\usepackage[T1]{fontenc}
\usepackage[french]{babel}
\usepackage[left=2cm,right=2cm,top=2cm,bottom=2cm]{geometry}
\usepackage{xcolor}
\usepackage{hyperref}
\usepackage{fontawesome5}
\usepackage{parskip}

\definecolor{primary}{RGB}{63, 150, 60}
\hypersetup{colorlinks=true, urlcolor=primary}

\begin{document}

\noindent
\begin{minipage}[t]{0.5\textwidth}
    \textbf{Loïc Aron MBASSI EWOLO}\\
    Élève-Ingénieur Bac+5 -- Double diplôme ENSEA\\[3pt]
    \faMapMarker\ Cergy, France\\
    \faPhone\ (+33) 7 59 52 33 98\\
    \faEnvelope\ \href{mailto:loicaron.mbassiewolo@ensea.fr}{loicaron.mbassiewolo@ensea.fr}
\end{minipage}
\hfill
\begin{minipage}[t]{0.4\textwidth}
    \raggedleft
    Cergy, le 10 mars 2026
\end{minipage}

\vspace{1.2cm}

\hfill
\begin{minipage}[t]{0.5\textwidth}
    \raggedleft
    \textbf{Schneider Electric}\\
    Service Recrutement\\
    Carros, France
\end{minipage}

\vspace{1cm}

\noindent\textbf{Objet :} Candidature en alternance -- Développeur logiciel embarqué réseaux industriels

\vspace{0.8cm}

Madame, Monsieur,

Intégrer un flux Lidar 3D sur UDP dans un système embarqué C/C++ sur Nvidia Jetson, synchroniser les données avec une caméra de profondeur et traiter le tout en temps réel pour de la cartographie autonome : c'est le travail que je mène dans le challenge CoHoMa (robotique militaire, ENSEA/Armée française). Cette expérience d'intégration de protocoles réseau dans un coupleur embarqué sous contrainte temps réel est directement transposable à l'intégration d'EtherCat dans le système Edge IO de Schneider Electric.

Ma programmation système C sous Linux (mémoire partagée, signaux POSIX, ordonnanceur FIFO) et mon prototypage IoT sur STM32 (capteurs IMU, fusion de données) m'ont donné une compréhension fine des couches basses : gestion mémoire, synchronisation, communication inter-processus. En parallèle, mon projet d'architecture microservices (RabbitMQ, API Gateway, Docker) m'a familiarisé avec l'intégration de protocoles de communication dans des systèmes existants et la mise en place de tests de vérification, compétences que je mettrai au service de l'intégration de la pile protocolaire EtherCat.

En stage à INRIA Saclay (équipe TRIBE, avril à août 2026), je développe des pipelines Python structurés avec intégration continue et documentation technique en anglais. Mon stage au laboratoire IoT de l'École Polytechnique de Yaoundé (TinyML, TF Lite, Arduino/Raspberry Pi) complète ce profil par l'optimisation de modèles pour microcontrôleurs sous contraintes mémoire strictes.

Je suis disponible dès septembre 2026 pour une alternance de 12 mois à Carros. Contribuer à l'intégration d'EtherCat dans un système d'entrées-sorties pour automates programmables, au sein d'une équipe spécialisée en logiciel embarqué, est un défi technique que je souhaite relever.

Je reste à votre disposition pour un entretien afin de discuter de ma candidature.

Je vous prie d'agréer, Madame, Monsieur, l'expression de mes salutations distinguées.

\vspace{0.8cm}

\textbf{Loïc Aron MBASSI EWOLO}\\
\textit{Élève-Ingénieur ENSEA / Polytechnique Yaoundé}

\end{document}
